\documentclass[10pt,a4paper]{amsart}
\usepackage[margin=19mm]{geometry}
\usepackage{amsmath,amssymb,mathtools}
\usepackage[scr=rsfs,cal=euler]{mathalfa}
\usepackage{xfrac}
\usepackage[unicode,hidelinks,bookmarks=false]{hyperref}
\newcommand{\ie}{i.e.\ }
\newcommand{\cf}{cf.\ }
\newcommand{\resp}{respectively\ }
\newcommand{\Subsection}{\subsection}
\providecommand{\nobreakdash}{\nobreak\hbox{-}}
\setlength{\emergencystretch}{2em}
\multlinegap0pt
\usepackage{float}
\usepackage{mathtools}
\usepackage{amssymb}
\usepackage[shortlabels]{enumitem}
\usepackage{xcolor}
\newtheorem{theorem}{Theorem}
\def\bigks{\mathbf{K}^s}
\def\gls#1{{\mathrm R}_{#1}}
\def\gsls#1{\gls{#1}^s}
\def\ls#1{\boldsymbol{\ell}(#1)}
\def\nn{\mathbb N}
\def\red#1{#1}
\def\rr{\mathbb R}
\def\valu#1{\mathfrak{s}(#1)}
\title{On the Maximization of Real Sequences}
\author{Assal\'e Adj\'e}
\date{Source: arXiv:2506.18409v2 (CC BY 4.0)}
\begin{document}
\maketitle

\noindent\emph{Source and licence.} On the Maximization of Real Sequences. Version arXiv:2506.18409v2 (CC BY 4.0), distributed by arXiv under the Creative Commons Attribution 4.0 International licence, http://creativecommons.org/licenses/by/4.0/, as recorded in the arXiv licence metadata for this exact version and shown on the abstract page https://arxiv.org/abs/2506.18409v2. Full licence notice: \texttt{licenses/arxiv-2506.18409-CC-BY-4.0.txt}.\par\smallskip

\noindent\textit{Editorial scope.} Counted unit: the article's theorem funatt (source0.tex lines 549--557). The article's complete original proof of it is delivered with it (source0.tex lines 559--573, one \texttt{proof} environment of 15 lines). No mathematics is written, repaired or re-derived: the statement and the proof are the article's own lines, in the article's own notation, and no retained line is substituted or re-flowed.  No further source-local context is needed: every cross-reference of every retained line resolves inside the two delivered spans.  Nothing was omitted from any retained span, so the package carries no omission record.  No retained line contains a citation, so the wrapper carries no bibliography and no bibliography entry.  No retained line contains a figure environment, an image-inclusion command or drawing code, so no image is delivered and the package declares no asset dependency.  Editorial formatting only, in addition: an AMS \texttt{amsart} preamble that restates the article's own declarations and macros used by the retained lines, namely the environments \texttt{theorem} on separate counters, together with the standard packages those lines need.  The retained environments print in this document as Theorem 1 in order of appearance, whereas the article prints the same items with the numbering of its own full document.  The complete native source of the article as distributed in the arXiv e-print for this exact version is bundled unchanged as \texttt{sources/180331/source0.tex}.
	\begin{theorem}[Existence of a useful optimal affine function] 
		\label{funatt}
		Let $u\in\Lambda$ such that $\gsls{u}\neq \emptyset$. Then, there exist $a>0$, $b\in (0,1)$ and $c\in\rr$ such that:
		\begin{enumerate}
			\item $\{k\in\nn : u_k>c\}\neq \emptyset$;
			\item $u_k\leq a b^k+c$ for all $k\in\nn$;
			\item $u_{\bigks_u}=ab^{\bigks_u}+c$.
		\end{enumerate}
	\end{theorem}

	\begin{proof}
		If $\gsls{u}$ is nonempty then $u_{\bigks_u}=\valu{u}>\ls{u}\in\rr\cup\{-\infty\}$ and for all $t\in (\ls{u},\valu{u})$ there exists 
		$k\in\nn$, such that $u_k>t$. So let $c$ be in $(\ls{u},\valu{u})$. Now as $c>\ls{u}$, there exists $N_c\in\nn$ s.t. $\sup_{k\geq N_c} u_k\leq c$. This is equivalent to $u_k\leq c$ for all $k\geq N_c$. We conclude that for all $k\geq N_c$, $u_k\leq c+ab^k$ for all $a>0$ and $b\in (0,1)$. Note that as $c<\valu{u}$ and $\bigks_u$ is the greatest element of the argmax of $u$, then $N_c$ is strictly greater than $\bigks_u$. Now, let us introduce the following set: 
		\[
		{AB}_c:= \left\{(a,b)\in\rr_+^* \times (0,1): u_{\bigks_u}=\valu{u}=a b^{\bigks_u}+c\right\}.
		\] 
		It is easy to see that ${AB}_c$ is nonempty, indeed, for all $b$ in $(0,1)$, $(\frac{u_{\bigks_u}-c}{b^{\bigks_u}},b)\in {AB}_c$. Moreover, if $\bigks_u>0$, for all $(a,b)\in {AB}_c$,
		for all $0<k< \bigks_u$, $u_k\leq u_{\bigks_u}=a b^{\bigks_u}+c<a b^k+c$ as $b\in (0,1)$. We also have $u_0\leq u_{\bigks_u}=ab^{\bigks_u}+c\leq a+c$ as $b\in (0,1)$. Hence, the result $u_k\leq ab^k+c$ holds on the sets $[0,\bigks_u]\cap \nn$ and $[N_c,+\infty)\cap\nn$ for every $(a,b)\in {AB}_c$. It suffices to find $(a,b)\in {AB}_c$, such that $u_k\leq ab^k+c$ for all $k\in (\bigks_u,N_c]\cap \nn$. Recall that we have $\bigks_u+1\leq N_c$.
		
		Let us write $(-\infty,0)\ni\gamma:=\max_{k\in (\bigks_u,N_c]\cap\nn} (u_{\bigks_u}-u_k)/(\bigks_u-k)$. This cannot be zero since $\bigks_u$ is the maximal integer $k$ such that $u_k=\max_{n\in\nn} u_n$. Now, let $(a,b)\in \rr_+\times (0,1)$ and define $f_{a,b}:\rr\ni y\mapsto a b^y+c$. The function $f_{a,b}$ is derivable, strictly convex and $f'_{a,b}(y)=ab^y\ln(b)$ for all $y\in\rr$. If $f'_{a,b}(\bigks_u)=\gamma$ and $(a,b)\in {AB}_c$ then, by the convexity of $f_{a,b}$, the tangent of $f$ at $y=\bigks_u$ satisfies for all $k\in\nn$, $f(k)\geq f'(\bigks_u)(k-\bigks_u)+f(\bigks_u)$ which is the same as $ab^k+c\geq \gamma (k-\bigks_u)+u_{\bigks_u}$. \red{From the} definition of $\gamma$, we have for all $k\in (\bigks_u,N_c\red{]\cap \nn}$, $(\bigks_u-k)\gamma\leq u_{\bigks_u}-u_k$ and so $u_k\leq u_{\bigks_u}+(k-\bigks_u)\gamma$. We conclude that if $f'_{a,b}(\bigks_u)=\gamma$ and $(a,b)\in {AB}_c$, $u_k\leq ab^k+c$ for all $k\in (\bigks_u,N_c]\red{\cap\nn}$. Finally, we look for $(a,b)\in \rr_+\times (0,1)$ such that:
		\[
		f_{a,b}(\bigks_u)=u_{\bigks_u} \text{ i.e. } (a,b)\in {AB}_c\text{ and } f'_{a,b}(\bigks_u)=\gamma.
		\]
		From $(a,b)\in {AB}_c$, we get $ab^{\bigks_u}=u_{\bigks_u}-c$. Hence, $f'_{a,b}(\bigks_u)=\gamma$ if and only if $\ln(b)=\frac{\gamma}{u_{\bigks_u}-c}<0$. So $b=\exp(\frac{\gamma}{u_{\bigks_u}-c})\in (0,1)$ and $a=(u_{\bigks_u}-c)\exp(-\frac{\bigks_u\gamma}{u_{\bigks_u}-c})>0$.
	\end{proof}

\end{document}
